\section{Sharpness}\label{sec:sharpness}

We finish the proof of Theorem~\ref{thm:main} by giving a direct
two-cycle construction. It also makes clear why a perturbative
calculation suffices for the lower bound, although it was not
used to prove the upper bound. The method is a direct first-order
return-displacement calculation, in the classical perturbative
tradition discussed in \cite{LMP,LlibreZhang}. We derive the exact
energy quotient needed for the argument below.

Set
\begin{equation}\label{eq:sharp-polynomial}
 F_\varepsilon(x)=\varepsilon f(x),\qquad
 f(x)=4x-\frac{20}{3}x^3+\frac85x^5.
\end{equation}
At $\varepsilon=0$, the solution through $(0,s)$, $s>0$, is
\[
 x(t)=s\sin t,\qquad y(t)=s\cos t.
\]
For $s$ in a neighborhood of either $1$ or $2$, smooth flow
dependence on a neighborhood of this compact circle, together
with transversality at $t=2\pi$, gives a smooth return time
$\tau(s,\varepsilon)$ near $2\pi$ and a smooth positive return
ordinate $P(s,\varepsilon)$ on $x=0,y>0$. These are the first
positive-axis returns for small $|\varepsilon|$. Indeed, away
from small time neighborhoods of the two axis crossings the
unperturbed $x$ is bounded away from zero; near each crossing
transversality gives a unique nearby crossing. These statements
hold on one common parameter neighborhood after shrinking the
two neighborhoods if necessary.

For $E=(x^2+y^2)/2$ the exact energy identity is
\[
 \dot E=-\varepsilon x f(x).
\]
Hence the one-turn energy displacement has the factorization
\begin{equation}\label{eq:energy-quotient}
 \frac{P(s,\varepsilon)^2-s^2}{2}
 =\varepsilon Q(s,\varepsilon),\qquad
 Q(s,\varepsilon)=
 -\int_0^{\tau(s,\varepsilon)}
 x(t;s,\varepsilon)f(x(t;s,\varepsilon))\,dt.
\end{equation}
The integral is smooth in $(s,\varepsilon)$ and defines the
quotient at $\varepsilon=0$ exactly. Using
\[
 \int_0^{2\pi}\sin^2t\,dt=\pi,\quad
 \int_0^{2\pi}\sin^4t\,dt=\frac{3\pi}{4},\quad
 \int_0^{2\pi}\sin^6t\,dt=\frac{5\pi}{8},
\]
we obtain
\begin{equation}\label{eq:sharp-averaged}
 Q(s,0)=-\pi s^2(4-5s^2+s^4).
\end{equation}
Its zeros at $s=1,2$ are simple:
\[
 Q_s(1,0)=6\pi,\qquad Q_s(2,0)=-48\pi.
\]
The implicit function theorem gives roots
$s_1(\varepsilon),s_2(\varepsilon)$ in disjoint positive
neighborhoods, for the same sufficiently small
$|\varepsilon|$. Fix any nonzero positive $\varepsilon$ in
that common range.

At these roots, $P+s>0$, so
Equation~\eqref{eq:energy-quotient} implies $P=s$.
Thus each root gives a periodic orbit. Differentiating the energy
displacement in $s$ at a root gives
\[
 s(P_s-1)=\varepsilon Q_s\ne0.
\]
Each return fixed point is therefore simple and isolated. The
two orbits are geometrically distinct, since their positive
intercepts lie in disjoint neighborhoods and a periodic orbit
has just one positive intercept, as proved in
Proposition~\ref{prop:cycle-matching}. After decreasing the common
upper bound on $\varepsilon>0$ if necessary, continuity from
$P_s(s,0)=1$ gives $P_s>0$ at both fixed points. The signs of $Q_s$
then give $P_s>1$ on the inner cycle and $0<P_s<1$ on the outer one.
Thus both cycles are hyperbolic: the inner one repels and the outer
one attracts. This proves sharpness.
\qed
